\begin{frame}{자주 하는 실수: $\rho$ 의 분산을 ``평균 에피소드
길이''에 대입하기}
  $\text{Var}[\rho] = 2^L - 1$ 은 \textbf{$L$ 이 같으면} 같은
  분산이지만, 실전에서는 \textbf{에피소드 길이가 고정되어
  있지 않다}. 5.2절 블랙잭처럼 ``언제 멈출지''가 변수인
  환경에서는 에피소드 길이가 2에서 100까지 \textbf{분포}.
  \vskip0.5em
  이 경우 $\text{Var}[\rho]$ 는 에피소드 \textbf{길이 분포}와
  행동/목표 정책의 차이 \textbf{모두}를 반영. ``에피소드 평균
  길이가 5니까 $\text{Var} \approx 31$ 이겠지''라고
  \textbf{평균 길이에 대입}하는 것은 \textbf{잘못}이다 ---
  \[
  \mathbb{E}[2^L] \ \ge\ 2^{\,\mathbb{E}[L]}
  \]
  (Jensen 부등식, $2^L$ 이 볼록) 이므로 \textbf{평균 길이에
  넣은 값보다 실제 분산이 더 크다}. 길이 분포가
  \textbf{넓을수록}(표준편차가 클수록) 격차는 커진다.
  \vskip0.4em
  \begin{block}{의미}
    이것이 ``에피소드 길이가 불규칙한 환경에서 off-policy 가
    \textbf{특히 어렵다}''는 말의 정확한 의미.
  \end{block}
\end{frame}
